\section{Formal Properties and Proof Sketches}
\label{app:proofs}

One short proof per Section~\ref{sec:evalprotocol} eval-metric primitive, matching the exact
parameterizations used there. \texttt{losses/}'s terms are not given formal proofs here ---
Section~\ref{sec:method}'s claims are validated empirically, via the unit tests in
Section~\ref{sec:results}/\texttt{tests/test\_losses\_sanity.py} and the ablation sweeps in
Sections~\ref{subsec:lossablation-setup}--\ref{subsec:lossablation-results}, rather than proven;
this is noted explicitly rather than left implicit.

\subsection{Shared Edge-Case Correctness}
\label{app:a1}

\textbf{Claim:} every primitive in Section~\ref{sec:evalprotocol} returns $1.0$ when both $P_c$
and $G_c$ are empty, and $0.0$ when exactly one is empty.

\textbf{Proof sketch:} verified by direct inspection of the identical guard pattern in each
implementation --- \texttt{dtaf1.py:43-51} ($n_{\mathrm{pred}}=0$ and $n_{\mathrm{gt}}=0 \to$ all
of precision/recall/f1 $=1.0$; either alone zero $\to$ all $0.0$), \texttt{cldice.py:51-56},
\texttt{boundary\_f1.py:68-72} (same shape, on skeleton and boundary pixel counts respectively),
\texttt{point\_f1.py:109-120} (on instance counts), and \texttt{apls.py:141-144} (on foreground
pixel counts, prior to any graph construction). Since every downstream composite (CBHM,
DTAF1-Topo, \texttt{cbhm\_soft}) is built from weighted means and harmonic means of these
primitives, this shared base case is what makes ``both classes vacuously correct'' a well-defined
$1.0$ at every level of the library, not just at the leaves.

\subsection{DTAF1 Boundedness and Reduction Equivalence}
\label{app:a2}

\textbf{Claim:} $\mathrm{dtaf1\_macro}$, $\mathrm{dtaf1\_weighted} \in [0,1]$ always; and
$\mathrm{dtaf1\_weighted} = \mathrm{dtaf1\_macro}$ when GT pixel counts are equal across all
classes in \texttt{class\_config}.

\textbf{Proof sketch:} each per-class $F1_c$ is itself a harmonic mean of two quantities in
$[0,1]$ (Section~\ref{subsec:dtaf1}), hence $F1_c \in [0,1]$; both the unweighted mean
(\texttt{dtaf1.py:124}) and the pixel-count-weighted mean (\texttt{dtaf1.py:125-127}) are convex
combinations of the $F1_c$ values, hence also in $[0,1]$. When every class's $n_c$ (GT pixel
count) is equal, the weights $n_c / \sum n_c$ reduce to $1/|C|$ for every class, making the
weighted sum numerically identical to the unweighted mean.

\subsection{CBHM Zero-Collapse}
\label{app:a3}

\textbf{Claim:} $\mathrm{cbhm}=0$ if and only if $\mathrm{clDice}_{\mathrm{mean}}=0$ or
$\mathrm{BF}_{\mathrm{mean}}=0$ (or both).

\textbf{Proof sketch:} directly from \texttt{unified.py:80-81} --- $\mathrm{score} =
2 \cdot \mathrm{cl\_mean} \cdot \mathrm{bf\_mean} / \mathrm{denom}$ if $\mathrm{denom}>0$ else
$0.0$. If either $\mathrm{cl\_mean}$ or $\mathrm{bf\_mean}$ is exactly 0, the numerator is 0 while
$\mathrm{denom}$ is the other (generally nonzero) term, so $\mathrm{score}=0$ exactly ---
regardless of how close to 1 the other mean is. Conversely, if both means are strictly positive,
both numerator and denominator are strictly positive, so $\mathrm{score}>0$. This is the formal
statement behind Table~\ref{tab:failure-modes}'s ``Road shifted 5px'' row ($\mathrm{clDice}=0$
forces $\mathrm{CBHM}=0$ even though $\mathrm{DTAF1}=1.000$ on the identical input).

\subsection{\texttt{cbhm\_soft} Non-Collapse}
\label{app:a4}

\textbf{Claim:} $\mathrm{cbhm\_soft}=0$ only if \emph{both} $\mathrm{cl\_mean\_weighted}=0$ and
$\mathrm{bf\_mean\_weighted}=0$ (unlike CBHM's single-sided collapse in
Section~\ref{app:a3}).

\textbf{Proof sketch:} from \texttt{unified.py:102}, $\mathrm{cbhm\_soft} = w_{\mathrm{cl}} \cdot
\mathrm{cl\_mean\_weighted} + w_{\mathrm{bf}} \cdot \mathrm{bf\_mean\_weighted}$ is a weighted
arithmetic mean with non-negative weights summing to 1 (or 0.5/0.5 when both totals are 0). A
weighted arithmetic mean of two non-negative terms is 0 only if every term with positive weight is
itself 0. Since $w_{\mathrm{cl}}, w_{\mathrm{bf}} \ge 0$ and at least one is strictly positive
whenever any GT pixels exist, $\mathrm{cbhm\_soft}=0$ requires whichever of
$\mathrm{cl\_mean\_weighted}$/$\mathrm{bf\_mean\_weighted}$ has positive weight to be 0 --- and if
only one type is present (the other has zero GT pixels and thus zero weight),
$\mathrm{cbhm\_soft}$ simply equals that one type's weighted mean, never forced toward 0 by the
\emph{other}, absent type. This is the mechanism behind the
\texttt{Khartoum\_img371}/\texttt{Khartoum\_img333} contrast in Section~\ref{subsec:weighted-reduction}:
\texttt{cbhm\_soft} can only be dragged down by a failing class in proportion to that class's own
pixel-count weight, not collapsed outright by it.

\subsection{APLS Score Boundedness}
\label{app:a5}

\textbf{Claim:} every per-pair APLS score, and hence the aggregate \texttt{apls} score, lies in
$[0,1]$.

\textbf{Proof sketch:} each per-pair score is either $0.0$ (unmatched, \texttt{apls.py:181}) or
$\max(0.0, 1.0 - |\mathrm{pred\_len} - \mathrm{gt\_len}|/\mathrm{gt\_len})$
(\texttt{apls.py:185}) --- the outer $\max(0.0, \cdot)$ clamps the lower bound, and the term
inside is $\le 1.0$ whenever $\mathrm{pred\_len} \ge 0$ (shortest-path lengths are non-negative by
construction), so every individual score lies in $[0,1]$. The aggregate \texttt{apls} score is
\texttt{np.mean(scores)}, a convex combination of values in $[0,1]$, hence itself in $[0,1]$.

\subsection{Point F1: Hungarian Optimality vs.\ Greedy Suboptimality}
\label{app:a6}

\textbf{Claim:} there exist configurations of GT/predicted point centroids where greedy
nearest-neighbor matching strictly under-counts true positives relative to the Hungarian
assignment used in \texttt{point\_f1} (\texttt{point\_f1.py:58-65}).

\textbf{Proof sketch (formalizing \texttt{TestPointInstanceMatching}'s adversarial case):}
consider two GT points $g_1, g_2$ and one predicted point $p$ positioned such that
$\mathrm{dist}(g_1,p) < \mathrm{dist}(g_2,p)$ but a \emph{second} predicted point $p'$ exists with
$\mathrm{dist}(g_2,p') < \mathrm{tolerance}$ while $\mathrm{dist}(g_1,p') > \mathrm{tolerance}$. A
greedy matcher processing predictions in an order that assigns $p \to g_1$ first (the globally
nearer pair) can leave $g_2$ unmatched to $p'$ if it was already ``claimed'' --- or, in the
canonical clustered-points case, greedy assigns the single nearest GT to \emph{both} predictions'
shared nearest neighbor, yielding $\mathrm{TP}=1$. The Hungarian assignment instead minimizes
total assignment cost over the \emph{complete} bipartite graph simultaneously
(\texttt{linear\_sum\_assignment} on the full cost matrix), correctly recovering the one-to-one
pairing $p{\leftrightarrow}g_1, p'{\leftrightarrow}g_2$ that gives $\mathrm{TP}=2$ --- strictly
more true positives from the same input, because it considers all pairings jointly rather than
greedily committing to the locally-best match first. This is also the formal reason
Section~\ref{subsec:heatmap-loss} does not attempt a differentiable relaxation of Hungarian
matching for the training loss: the property that makes it correct (global joint optimization
over the complete bipartite graph) is exactly the discrete, combinatorial structure that has no
useful gradient.
